\section{Curvature homotopies for actual triples}
\label{sec:r10-triple-curvature-homotopy}

Fix the actual pair and triple spaces of CPT1.  In degree $s$ put
\[
 L_3^s=
 B^{s-3,s}\oplus B^{s,s-3}
       \oplus B^{s-2,s-1}\oplus B^{s-1,s-2},
 \qquad
 \partial_3(U,V,P,Q)=
       (d'U+d''P,\ d''V+d'Q,\ d'P+d''Q).
 \tag{QH1}
\]
The target order is $(A,B,M)$ as in CPT1.  Negative degrees mean
zero.  Full anticommutation gives $\partial_2\partial_3=0$.

Let $h$ be any actual diagonal form of degree $(k,k)$, $k\ge0$,
and $q=(A,B,M)\in L_2^s$.  Define
\[
 K_h(q)=
       (d''h\wedge A,\ -d'h\wedge B,\
                         -h\wedge d''M,\ -d'h\wedge M)
                         \in L_3^{s+k+1}.
 \tag{QH2}
\]
The frozen CPT homotopy for a pair is
$H_h(p)=(d''h\wedge p',-d'h\wedge p'',h\wedge d'p')$.

\begin{theorem}[The full nonclosed-triple identity]
For every actual triple, without a closedness hypothesis,
\[
             \partial_3K_h(q)=Dh\wedge q-H_h(\partial_2q).
 \tag{QH3}
\]
In particular, when $\partial_2q=0$, multiplication by $Dh$
is an explicitly displayed $\partial_3$ boundary.
\end{theorem}

\begin{proof}
Write $p=\partial_2q$, so
$p'=d'A+d''M$ and $p''=d''B+d'M$.  The total degree of
$h,A,B,M$ is even.  Applying both full Leibniz rules, the first
component of the left side of (QH3) is
\[
 d'(d''h\wedge A)+d''(-h\wedge d''M)
       =Dh\wedge A-d''h\wedge(d'A+d''M)
       =Dh\wedge A-d''h\wedge p'.
 \tag{QH4}
\]
The second is
\[
 d''(-d'h\wedge B)+d'(-d'h\wedge M)
       =Dh\wedge B+d'h\wedge(d''B+d'M)
       =Dh\wedge B+d'h\wedge p''.
 \tag{QH5}
\]
The third is
\[
 d'(-h\wedge d''M)+d''(-d'h\wedge M)
               =Dh\wedge M-h\wedge DM.
 \tag{QH6}
\]
Since $d'p'=DM$, these are exactly the three components on
the right of (QH3).  No closedness or cancellation of a curvature
factor was used.
\end{proof}

\subsection{A finite transport criterion at the next cochain level}

Let $W=N\setminus C$, and suppose $R\in B^{1,1}(N)$ is closed
under both full differentials.  Let an actual triple
$e\in L_2^a(N)$ satisfy $\partial_2e=0$ on all of $N$.
Let $C_j=R-D\Phi_j$ on $W$, where the
$\Phi_j$ are actual scalars, and let $P$ be a homogeneous real
polynomial of degree $d$.  Use the explicit finite CPT7--CPT9
potential $h_P$:
\[
             Dh_P=P(C_1,\ldots,C_t)-P(1,\ldots,1)R^d.
 \tag{QH7}
\]
For $d\ge1$, (QH3) gives the literal actual identity
\[
 P(C_1,\ldots,C_t)e
        =\partial_3K_{h_P}(e)
                     +P(1,\ldots,1)R^d e.
 \tag{QH8}
\]
For $d=0$ the equality is immediate with zero quadruple.  The
last triple in (QH8) is actual on the unpunctured $N$.  The
coefficient triple $e$ must be independently $\partial_2$-closed;
closedness of a product $P(C)e$ is not sufficient.

Thus a separately verified actual factorization
\[
            q_1-q_0=\partial_3t+
                   \sum_jP_j(C_{j,1},\ldots,C_{j,t_j})e_j,
 \qquad \partial_2e_j=0\ \hbox{on }N,
 \tag{QH9}
\]
gives an explicit quadruple-plus-extending-triple witness by
adding the $K_{h_{P_j}}(e_j)$ to $t$.  The equality in (QH9)
is a substantive supplied identity.  This theorem does not
supply it for a general triple-overlap or sorting cube.

\subsection{The exact three-swap expression still to be analysed}

For three distinct diagonal factors $A,B,C$ of weights $a,b,c$,
write $q_{a,b}(f)$ for the specific capped-square triple CS36
computed on the indicated intermediate tuple, with the same
fixed positive cap constant.  The two adjacent-swap paths from
$ABC$ to $CBA$ give the actual triple difference
\[
 \begin{split}
 \mathcal Q_{a,b,c}(f)={}&q_{a,b}(f)+q_{a,c}(Bf)+q_{b,c}(f)\\
                      &-q_{b,c}(Af)-q_{a,c}(f)-q_{a,b}(Cf).
 \end{split}
 \tag{QH10}
\]
The primitive endpoints agree, and the extending pairs on the
two paths are literally the same sum
$(e_{a,b}+e_{a,c}+e_{b,c})R^{r-2}$.  Therefore CS3 gives
\[
                         \partial_2\mathcal Q_{a,b,c}(f)=0.
 \tag{QH11}
\]
This conclusion is only closedness of a specified actual triple.
It is not a $\partial_3$ witness.

A simple independently closed extending-triple family is
available for potential comparisons: take the PPC tetrahedron
triple on four pluriharmonic scalars on $N$.  Its tetrahedron
boundary vanishes because all their curvatures relative to $R$
are equal.  Such triples can therefore be used in (QH8) when
an actual factorization involving them has been proved.  This
observation does not identify the concrete cube (QH10) with
that family.  Its different square terms use different passive
blocks and different capped norm potentials; those differences
must be compared by actual cochain identities.

The remaining exact task is to construct
$\mathcal Q_{a,b,c}=\partial_3t+e|_W$ with a displayed actual
quadruple $t$ and triple $e$ defined across cancellation, or to
give a precisely scoped obstruction for these chosen square
triples.  Neither QH3 nor the five-potential PPC identity
automatically supplies this factorization.  It is consequently
not inferred that arbitrary multi-open frame paths glue.

All operations here are actual finite differential and wedge
operations, commute with restriction and refinement, and retain
all mixed coefficient jets on a fixed whole presentation.
The spaces have their own inherited actual coefficient-final
topologies.  No Hausdorff quotient, whole-kernel comparison,
general growth bound, higher-form descent or BB-height conclusion
is asserted.  The original field and absolute value are unchanged.

\paragraph{Status and construction.}
QH3 and QH8 are R10 finite actual identities constructed by the
diagonal constructor.  QH11 is a direct consequence of the
audited CS square theorem.  The explicit cube's relative
$\partial_3$ factorization is OPEN here.  An independent reviewer
must be distinct from the constructor and bind this exact source.

